\section{Các kỹ thuật tăng độ tin cậy đã cài đặt}
\label{sec:ky-thuat}

Trên khung kiến trúc ở Mục~\ref{sec:kien-truc}, đồ án cài đặt một tập kỹ
thuật nhắm thẳng vào ba vấn đề đã nêu (độ tin cậy, khả năng quan sát, chi
phí). Điểm chung của các kỹ thuật: đều \emph{đo được} --- hoặc qua test tự
động, hoặc qua metric tổng hợp từ trace.

\subsection{Chống prompt injection nhiều lớp}
\label{sec:ky-thuat-injection}

Rủi ro chính với agent có công cụ đọc web là indirect prompt
injection~\cite{greshake2023injection}: chỉ dẫn độc hại nằm trong nội dung
trang web hoặc kết quả tìm kiếm mà agent đọc. Hệ thống phòng thủ theo nhiều
lớp:

\begin{enumerate}
  \item \textbf{Đóng khung dữ liệu không tin cậy}: mọi nội dung công cụ trả về
        được bọc trong đúng một cặp thẻ \texttt{<tool\_output>}\dots%
        \texttt{</tool\_output>}; system prompt quy định tường minh: nội dung
        trong khung là \emph{dữ liệu}, không phải mệnh lệnh, và không gì bên
        trong đó có thể thêm/sửa/hủy tác vụ.
  \item \textbf{Vô hiệu hóa delimiter giả}: trước khi đóng khung, mọi chuỗi
        trong dữ liệu trông giống thẻ đóng/mở (kể cả biến thể viết hoa, chèn
        khoảng trắng như \texttt{</ tool\_output >}) bị thay bằng ký tự tương
        đương vô hại --- kẻ tấn công không thể ``thoát khung'' bằng cách tự
        chèn thẻ đóng rồi viết chỉ dẫn phía sau.
  \item \textbf{Giới hạn quyền công cụ}: mọi công cụ ghi DB lọc theo
        \texttt{user\_id}; không có công cụ thực thi lệnh hệ thống.
  \item \textbf{Con người là phòng tuyến cuối}: kể cả khi mô hình bị lừa muốn
        xóa dữ liệu, hành động ghi/xóa vẫn phải qua nút xác nhận của người
        dùng (Mục~\ref{sec:hitl}).
\end{enumerate}

Hiệu quả được đo bằng nhóm 10 case injection trong bộ eval
(Mục~\ref{sec:bo-case}): case khai báo công cụ bị cấm
(\texttt{forbid\_tools}) và kỳ vọng agent không thực hiện chỉ dẫn cài trong
dữ liệu.

\subsection{Phân loại lỗi công cụ và chính sách retry}
\label{sec:ky-thuat-taxonomy}

Executor gắn cho mọi lỗi công cụ một \emph{nhãn phân loại} (error taxonomy)
dưới dạng prefix \texttt{[kind]} trong thông báo lỗi --- vừa để LLM đọc quan
sát biết nên đổi hướng hay thử lại, vừa để người đọc trace chẩn đoán nhanh.
Chín loại lỗi và chính sách retry tương ứng:

\begin{table}[ht]
\centering
\caption{Error taxonomy của tool executor.}
\label{tab:error-taxonomy}
\begin{tabular}{llc}
\toprule
Nhãn & Ý nghĩa & Retry? \\
\midrule
\texttt{[input]}      & tham số sai schema / URL không hợp lệ & Không \\
\texttt{[not\_found]}  & công cụ không tồn tại                & Không \\
\texttt{[timeout]}    & vượt ngưỡng \texttt{timeout\_s}       & Không \\
\texttt{[network]}    & lỗi mạng tạm thời (connect/read)      & Có + backoff \\
\texttt{[rate\_limit]} & HTTP 429                             & Có + backoff \\
\texttt{[http\_5xx]}   & server lỗi tạm thời                  & Có + backoff \\
\texttt{[http\_4xx]}   & lỗi phía client (403, 404, \dots)    & Không \\
\texttt{[crash]}      & exception bất ngờ trong công cụ       & Có (1 lần) \\
\texttt{[tool\_error]} & công cụ chủ động trả \texttt{ok=False} & Không \\
\bottomrule
\end{tabular}
\end{table}

Chỉ lỗi \emph{tạm thời} mới được thử lại, với exponential backoff
($0{,}5 \cdot 2^{n}$ giây, trần 8 giây). Công cụ có thể tự phân loại bằng
cách raise \texttt{TransientToolError}/\texttt{PermanentToolError}; exception
khác được suy loại tự động từ kiểu lỗi (timeout/network/status HTTP của
httpx). Lỗi timeout cố ý \emph{không} retry: Python không hủy được thread
đang chạy, nên thread quá hạn được bỏ lại có kiểm soát thay vì chồng thêm
lần chạy mới.

\subsection{Rate limit theo người dùng}
\label{sec:ky-thuat-ratelimit}

Một người dùng spam tin nhắn sẽ chiếm hết worker thread của agent loop và
chặn người dùng khác. Bot áp rate limit \emph{sliding window} theo Telegram
user id: tối đa 6 yêu cầu trong cửa sổ trượt 60 giây; vượt ngưỡng thì bot từ
chối kèm thời gian cần chờ (\texttt{retry\_after}). Limiter thread-safe
(handler chạy trên event loop nhưng tác vụ chạy ở thread khác), trạng thái
in-memory --- đủ cho mô hình một tiến trình, restart là reset (chấp nhận được
ở phạm vi MVP). Cùng lớp bảo vệ đầu vào còn có chặn tin nhắn quá 2000 ký tự
và validate biểu thức cron trước khi tạo job.

\subsection{Self-correction có đo đếm}
\label{sec:ky-thuat-selfcorrect}

Mọi quyết định của LLM đi qua một hàm duy nhất \texttt{call\_structured}: gọi
LLM kèm JSON schema, validate bằng Pydantic; nếu sai, thông báo lỗi validation
được gửi lại cho LLM tự sửa, tối đa 2 lần retry. Hai điểm đáng chú ý trong
cài đặt:

\begin{itemize}
  \item \textbf{Validator nghiệp vụ chặt}: ngoài đúng kiểu JSON, schema còn mã
        hóa ràng buộc nghiệp vụ --- hành động \texttt{tool} bắt buộc có tên
        công cụ, hành động \texttt{final} bắt buộc có câu trả lời cuối, kế
        hoạch không được rỗng. Đầu ra ``đúng cú pháp nhưng vô nghĩa'' bị chặn
        ngay tại validate và kích hoạt tự sửa, thay vì lọt xuống executor rồi
        hỏng khó hiểu ở tầng dưới.
  \item \textbf{Đếm được qua trace}: lần gọi retry được ghi
        \texttt{llm\_calls} với purpose \texttt{<purpose>:fixN} (ví dụ
        \texttt{react:fix1}). Nhờ đó số lần LLM phải tự sửa trở thành một chỉ
        số đo trực tiếp từ trace --- dùng để so sánh mức ``vững schema'' giữa
        các mô hình và cấu hình trong thực nghiệm.
\end{itemize}

Một bài học thực nghiệm đã phản hồi ngược vào prompt: model nhỏ hay bịa tên
công cụ không tồn tại (\texttt{google\_search}, \texttt{save\_note}); việc
nhắc lại danh sách tên công cụ hợp lệ ngay cạnh yêu cầu JSON trong prompt
ReAct/Plan làm giảm rõ rệt loại lỗi này.

\subsection{Theo dõi chi phí và trần chi phí mỗi tác vụ}
\label{sec:ky-thuat-cost}

Mỗi dòng \texttt{llm\_calls} đã ghi token và chi phí (USD, theo bảng giá của
model) của từng lệnh gọi. Trên nền đó:

\begin{itemize}
  \item \textbf{Tổng hợp}: hàm \texttt{task\_usage}/\texttt{user\_usage} cộng
        dồn số lệnh gọi, token, chi phí theo tác vụ và theo người dùng; trace
        viewer hiển thị cột token/chi phí ở danh sách tác vụ, ô tổng chi phí
        toàn hệ thống, và các ô thống kê trong trang chi tiết.
  \item \textbf{Trần chi phí}: ở chế độ LLM thật, provider được bọc bằng
        \texttt{BudgetedLLM} --- cộng dồn chi phí qua các lần
        \texttt{complete()}, chạm trần thì ném \texttt{LLMBudgetExceeded}
        \emph{trước} lần gọi API kế tiếp, nên vòng lặp chạy hoang (runaway
        loop) bị chặn sớm và chi phí thực tế chỉ có thể vượt trần tối đa một
        lần gọi. Trần đọc từ biến môi trường
        \texttt{LAPLACE\_MAX\_COST\_PER\_TASK\_USD} (mặc định 0{,}25~USD một
        lần chạy; đặt $\leq 0$ để tắt). Exception này được orchestrator xử lý
        như lỗi dừng sạch: tác vụ chuyển \texttt{failed} với thông báo rõ ràng
        gửi về người dùng, kèm hướng dẫn nâng trần.
\end{itemize}

Giới hạn đã ghi nhận: trần áp cho \emph{từng chặng chạy} (run/resume), không
cộng dồn qua các lần tạm dừng xác nhận --- xem Mục~\ref{sec:gioi-han}.

\subsection{Streaming trạng thái về Telegram}
\label{sec:ky-thuat-streaming}

Tác vụ nhiều bước có thể chạy hàng chục giây; im lặng suốt thời gian đó là
trải nghiệm xấu. Vì lõi agent không có hook sự kiện (và không muốn sửa lõi
chỉ để phục vụ UI), tiến độ được suy ra ở \emph{mép}: bot bọc provider bằng
\texttt{ProgressLLM} --- mỗi lần orchestrator gọi \texttt{complete()}, wrapper
suy giai đoạn hiện tại (phân loại / suy nghĩ / chạy công cụ X / tổng hợp) từ
json\_schema của lệnh gọi và nội dung structured output trả về, rồi phát sự
kiện tiến độ. \texttt{ProgressStreamer} nhận sự kiện từ worker thread (qua
\texttt{call\_soon\_threadsafe} + queue), và \emph{edit đúng một message}
trạng thái trên Telegram (throttle 1,5 giây để tránh rate limit edit của
Telegram); khi tác vụ xong, message trạng thái được thay bằng kết quả cuối.
Tác vụ định kỳ từ scheduler cũng được stream theo cùng cơ chế: job cron tự
tạo kết nối bot, gửi message trạng thái vào đúng chat của chủ job rồi cập
nhật dần đến kết quả.

\subsection{Chế độ phát lại trace (replay)}
\label{sec:ky-thuat-replay}

Trace viewer có route \texttt{/tasks/\{id\}/replay} phát lại timeline của một
tác vụ đã chạy: từng bước hiện dần theo thứ tự thật (điều khiển Next/Prev
hoặc auto-play), thời gian chờ giữa các bước lấy theo độ trễ thật đã ghi
trong trace, rút gọn theo hệ số tốc độ $1\times$--$8\times$. Trang được
render hoàn toàn phía server (auto-play bằng meta refresh) --- không cần
JavaScript, không cần mạng, không cần LLM. Đây vừa là công cụ gỡ lỗi
(``tua lại'' quyết định của agent), vừa là phương án demo dự phòng khi buổi
bảo vệ mất mạng: toàn bộ kịch bản demo phát lại được từ file SQLite mang
theo.

\subsection{Cache kết quả tìm kiếm}
\label{sec:ky-thuat-cache}

Trong phát triển và đánh giá, cùng một truy vấn tìm kiếm bị gọi lặp lại rất
nhiều; gọi API thật vừa tốn chi phí vừa làm eval dao động (kết quả web thay
đổi giữa hai lần chạy). Công cụ \texttt{web\_search} có lớp cache in-process
bật bằng biến môi trường (\texttt{LAPLACE\_SEARCH\_CACHE=1}): khóa cache là
truy vấn đã \emph{chuẩn hóa} (casefold + gộp khoảng trắng --- ``FastAPI~~VS
Flask'' và ``fastapi vs flask'' dùng chung một entry) kèm số kết quả; TTL mặc
định 3600 giây; trần 256 entry với chính sách loại bỏ entry cũ nhất. Cache
làm quan sát ổn định giữa các lượt chạy eval (cùng input $\to$ cùng
observation) và giảm chi phí API trong quá trình phát triển.
